\section{Rutas de razonamiento para la automejora}
\subsection{De Chain-of-Thought a la autorrecompensa}
Weston subraya que el crecimiento de la capacidad de razonamiento no puede depender solo de un prompt fuerte, sino que requiere que el propio modelo, durante la etapa de entrenamiento, «piense sobre su propio pensamiento». De ahí surgen los Self-Rewarding language models: durante el entrenamiento, el modelo desempeña a la vez el papel de generador de problemas, de quien los resuelve y de evaluador, y sustituye la calificación humana por una autoevaluación, de manera similar a como el mecanismo de ``verified response'' del Self-Feeding Chatbot se extiende a las tareas de razonamiento. Este proceso busca eliminar la confusión entre «memoria» y «razonamiento», por lo que los base data siguen incluyendo ejemplos que fuerzan el chain-of-thought, de modo que el reward model pueda identificar qué parte del reasoning explica la respuesta.

\subsection{El equilibrio entre el razonamiento y la autoevaluación}
Durante las iteraciones de self-improvement, el modelo cae fácilmente en un razonamiento encadenado excesivo: para obtener un reward más alto, genera chains cada vez más largas pero no necesariamente más correctas. Weston introduce el contraste ``self-rewarding vs self-feeding'' y señala que este último depende de la retroalimentación conversacional, mientras que el primero depende de la autoverificación (self-verification). Propone insertar ``cross checking'' en el entrenamiento self-rewarding: si el modelo no logra rechazar un error en el draft inicial, se le hace ``preguntarse a sí mismo'' una vez más.

\begin{warningbox}{Evitar el overthinking}
Por razones de seguridad, es necesario limitar la expansión del Chain-of-Thought, pues de lo contrario el inference compute se desperdicia en pensamientos repetidos. La introducción del Meta judge responde en parte a la necesidad de detectar a tiempo los casos en que ``incluso con una cadena de razonamiento larga, la respuesta sigue siendo incorrecta''.
\end{warningbox}

\subsection{Señales de entrenamiento y modelos de recompensa}
En la clase, Weston usa tres conceptos de reward para explicar el origen de la señal: el Process Reward Model (PRM, recompensa paso a paso), el Outcome Reward Model (ORM, que solo observa la respuesta final) y el RLHF / DPO tradicional (preferencia humana). El entrenamiento self-rewarding se acerca más al PRM, pues intenta estimar una puntuación después de cada paso de razonamiento del modelo, mientras que el Meta judge adopta la perspectiva del ORM para comparar el draft con la verified response. Weston recurre a ejemplos de benchmark como ``DeepSeek-R1, O1-R1'' para mostrar que, sin evaluación humana adicional, la autorrecompensa aún puede elevar el reasoning win rate (sobre todo en Alpaca Eval 2).

\begin{importantbox}{Comparación de modelos de recompensa}
\begin{itemize}
    \item PRM (recompensa de proceso): cada paso de la chain tiene señal, adecuado para ajustar el reasoning trace.
    \item ORM (recompensa de resultado): solo observa el juicio final, adecuado para casos como el meta judge, que solo necesita determinar cuál respuesta es mejor.
    \item RLHF / DPO: requiere preferencia humana, el costo de entrenamiento es alto, y aun así se sigue usando como baseline para la comparación.
\end{itemize}
\end{importantbox}

\begin{knowledgebox}{El valor de la recompensa verificable}
Weston señala la idea de ``verifiable reward'': si el modelo puede tratar su propia verified response como un ground truth aproximado, entonces se pueden usar estos data para train el reward model y reforzar la self-evaluation, en lugar de depender por completo de etiquetas humanas de bueno o malo.
\end{knowledgebox}

\subsection{Resumen de la sección}
El núcleo de la automejora consiste en reescribir la reward signal: dejar que el propio modelo determine el valor de la cadena de razonamiento, y conservar durante el entrenamiento tanto la process view como la outcome view. Los capítulos siguientes concretarán este objetivo de entrenamiento en un pipeline reutilizable.

